\documentclass[10pt]{amsart}
\usepackage[a4paper,margin=22mm]{geometry}
\usepackage{fontspec}
\setmainfont{Latin Modern Roman}[SmallCapsFont={lmromancaps10-regular.otf}]
\usepackage{amsmath,amssymb,amsthm,mathtools,hyperref,xcolor}
\usepackage[shortlabels]{enumitem}
\usepackage[numbers,square]{natbib}
\newtheorem{theo}{Theorem}[section]
\newtheorem{lemm}[theo]{Lemma}
\newtheorem{Proposition}[theo]{Proposition}
\theoremstyle{definition}
\newtheorem{Definition}[theo]{Definition}
\newtheorem{exam}[theo]{Example}
\newcounter{cdrthm}
\newenvironment{enonce}[2][]{\refstepcounter{cdrthm}\par\smallskip\noindent\textbf{#2~\thecdrthm}.\ }{\par\smallskip}
\newcommand{\ie}{i.e.\ }
\newcommand{\eg}{e.g.\ }
\newcommand{\resp}{respectively\ }
\emergencystretch=3em

\usepackage{tikz, float} \usetikzlibrary {positioning}

\newcommand{\BLambda}{{B\Lambda}}

\DeclareMathOperator{\Gr}{Gr}
\DeclareMathOperator{\id}{id}
\DeclareMathOperator{\ini}{in}
\DeclareMathOperator{\vol}{vol}
\DeclareMathOperator{\sgn}{sgn}



\newcommand{\I}{\mathcal{I}}
\newcommand{\K}{\mathbb{K}}
\newcommand{\TP}{\mathbb{TP}}


\tikzstyle{Cgray}=[scale = .45,circle, fill = gray, minimum size=3mm] \tikzstyle{Cblack}=[scale = .7,circle, fill = black, minimum size=3mm]
%\tikzstyle{Cred}=[scale = .45,circle, fill = red, minimum size=3mm]
\tikzstyle{Cblue}=[scale = .5,circle, fill = blue, inner sep = 0pt, minimum
size=3mm]
\tikzstyle{C1}=[scale = .7,circle, fill = black!0, inner sep = 0pt, minimum
size=3mm]

\tikzstyle{test2}=[scale = 1.5,circle, fill = black!0, inner sep = 0pt, minimum
size=3mm]
\tikzstyle{Cwhite}=[scale = .45,circle, fill = white, minimum size=3mm] 
%\tikzstyle{Cgray}=[scale = .3,circle, fill = gray!20, minimum size=3mm] 
\tikzstyle{Cblack2}=[scale = .3,circle, fill = black, minimum size=3mm] 
\tikzstyle{Cblack}=[scale = .3,circle, fill = black, minimum size=3mm]
\tikzstyle{C0}=[scale = .9,circle, fill = black!0, inner sep = 0pt, minimum size=3mm]
\tikzstyle{C1}=[scale = .7,circle, fill = black!0, inner sep = 0pt, minimum size=3mm]
\tikzstyle{Cred}=[scale = .4,circle, fill = red, minimum size=3mm] 
\tikzstyle{Cblack3}=[scale = .4,circle, fill = black, minimum size=3mm] 


%%%%%%%%% a placer avant \begin{document}


%%%%%%%%%%%%%%%%%%%%%%%%%%%%%%%%%%%%%
\newcommand*{\mk}{\mkern -1mu}
\newcommand*{\Mk}{\mkern -2mu}
\newcommand*{\mK}{\mkern 1mu}
\newcommand*{\MK}{\mkern 2mu}

\hypersetup{urlcolor=purple, linkcolor=blue, citecolor=red}


\newcommand*{\romanenumi}{\renewcommand*{\theenumi}{\roman{enumi}}}
\newcommand*{\Romanenumi}{\renewcommand*{\theenumi}{\Roman{enumi}}}
\newcommand*{\alphenumi}{\renewcommand*{\theenumi}{\alph{enumi}}}
\newcommand*{\Alphenumi}{\renewcommand*{\theenumi}{\Alph{enumi}}}
%%%%%%%%%%%%%%%%%%%%%%%%%%%%%%%%%%

\title[Quadratic two-block matching-field ideals]{Quadratic generation of three-row two-block diagonal matching-field toric ideals}
\author{Fatemeh Mohammadi; Kristin Shaw}
\date{}
\begin{document}
\maketitle
\noindent\textbf{Source.} Original Theorem 1.4, Fatemeh Mohammadi and Kristin Shaw, Toric degenerations of Grassmannians from matching fields, Algebraic Combinatorics 2 (2019), 1109--1124, DOI 10.5802/alco.77. The original theorem, necessary tableau/kernel/block definitions and four-case proof with its original worked examples follow. Native source and publication are preserved.
\setcounter{section}{1}
\setcounter{theo}{3}
\begin{theo}\label{Theorem:2block}
The ideal of a $2$-block diagonal matching field of size $3 \times n$ is quadratically generated. 
\end{theo}
\section{Preliminaries}\label{prelim} 

Throughout we set $[n]\coloneqq \{1, \ldots , n\}$ and we use $\mathbf{I}_{k, n}$ to denote the collection of subsets of $[n]$ of size $k$. The symmetric group on $k$ elements is denoted by $S_k$. We also fix a field $\mathbb{K}$. 

\smallskip

The Grassmannian $\Gr(k, n)$ is the space of all $k$ dimensional linear subspaces of $\mathbb{K}^n$. A point in $\Gr(k, n)$ can be represented by a $k\times n$ matrix with entries in $\mathbb{K}$. 
Let $X=(x_{ij})$ be a $k\times n$ matrix of indeterminates. For a subset $I = \{i_1,\ldots,i_k\} \in \mathbf{I}_{k, n}$, let $X_I$ denote the $k\times k$ submatrix with the column indices $i_1,\ldots,i_k$. 
The \emph{Pl{\"u}cker forms (or Pl{\"u}cker coordinates)} are $P_I = \det (X_I)$ for $I \in\mathbf{I}_{k, n}$. These forms determine the \emph{Pl{\"u}cker embedding} from
$ \Gr (k,n)$ into $\mathbb{P}^{\binom{n}{k}-1}$. 

In the following, we consider the polynomial ring $\mathbb{K}[x_{ij}]$ on the variables $x_{ij}$ with $1\leq i\leq k$ and $1\leq j\leq n$ and the polynomial ring $\mathbb{K}[P_I]$ on the Pl{\"u}cker variables with $|I|=k$.

\begin{Definition}\label{def:ideal}The Pl{\"u}cker ideal $\I_{k,n}$ is defined as the kernel of the map 
\begin{equation}\label{eqn:pluckermap}
\begin{aligned} 
\psi \colon \mathbb{K}[P_I] &\rightarrow \mathbb{K}[x_{ij}] 
\\ 
 P_{I} &\mapsto \det (X_I).
\end{aligned}
\end{equation}
The \emph{Pl{\"u}cker algebra} is the finitely generated algebra $\mathbb{K}[P_I ]/{\mathcal{I}_{k,n}}$ 
denoted by $\mathcal{A}_{k,n}$.
\end{Definition}




\begin{Definition}\label{def:matching}
A $k\times n$ matching field is a map $\Lambda: \mathbf{I}_{k,n} \to S_k$.
\end{Definition}

Given a $k\times n$ matching field $\Lambda$ and a subset $I = \{i_1, \ldots , i_k\} \in \mathbf{I}_{k, n}$ 
we consider the set to be ordered by $i_1 < \cdots < i_{k}$. We think of the permutation {$\sigma=\Lambda(I)$} as inducing a new ordering on the elements of $I$ 
where the position of $i_s$ is determined by the value of $\sigma(s)$.
It is convenient to represent the variable $P_I$ as a $k \times 1$ tableau 
 where $(\sigma(r), 1)$ contains $i_{r}$.

\begin{Definition}
Let $\Lambda$ be a size $k \times n$ matching field. 
A $\Lambda$-tableau is a tableau of size $k \times d$ for any $d\geq 1$ with entries in $[n]$,
so that the entries in each column are pairwise distinct and filled according to the order determined by $\Lambda$. 
\end{Definition}


\begin{exam}\label{ex:diagonal1}
The \emph{diagonal matching field} assigns to
each subset $I \in \mathbf{I}_{k,n}$ the identity permutation~\cite[Example~1.3]{sturmfels1993maximal}.
Therefore, a $\Lambda$-tableau is a rectangular tableau of size $k \times d$ filled with entries in $[n]$ such that the columns are strictly increasing. 
\end{exam}

\begin{exam}\label{ex:pointed1}
A $k \times n$ matching field is called \emph{pointed} if there exists $i_1, \ldots, i_k \in [n]$ such that if $i_s \in I$ for some 
$1 \leq s \leq k$ then $\Lambda(I)(s) = s$~\cite[Example~1.4]{sturmfels1993maximal}. In other words, if $i_s \in I$ for some $1\leq s \leq k$ then the column corresponding to $P_I$ contains $i_s$ in row $s$. Below are the tableaux representing $P_I$ for a pointed matching field of size $3 \times 5$ which is otherwise filled diagonally:
\[
\begin{array}{c}1 \\2 \\ 3 \end{array} , \quad
\begin{array}{c} 1 \\2 \\4 \end{array},\quad
\begin{array}{c}1 \\2 \\ 5 \end{array} , \quad
\begin{array}{c}1 \\ 4 \\ 3 \end{array} ,\quad
\begin{array}{c}1 \\5 \\ 3\end{array} ,\quad
\begin{array}{c}4 \\ 2\\ 3 \end{array} ,\quad
\begin{array}{c}5 \\2 \\ 3\end{array} ,\quad
\begin{array}{c} 4 \\ 2 \\ 5 \end{array}. 
\]
Generalising this, we say that a size $k \times n$ matching field $\Lambda$ is pointed on $S \subset [n]$ if for all $i \in S$ there exists a $j_i$ such that $i$ always appears in row $j_i$ of any $\Lambda$ matching field tableau. 
Here $S$ need not have size equal to $k$. 
\end{exam}


A monomial $\Pi_{I \in A} P_I $ can be represented by a $\Lambda$-tableau of size $k \times |A|$ given by the concatenation of the columns with content $I$ filled according to the matching field. 
To each monomial $\Pi_{I \in A} P_I $ we associate a sign $\epsilon_{A} = \pm 1$ determined by the signature of the permutations $\Lambda(I)$ for all $I \in A$. More precisely, 
\[
\epsilon_{A} = \Pi_{I \in A} \sgn (\Lambda(I)).
\]


\begin{Definition}
Given a matching field $\Lambda$, the \emph{matching field ideal} $\mathcal{I}_{\Lambda} \subset \mathbb{K}[x_{ij}]$ is generated by the 
 binomial relations 
\begin{equation}\label{eq:binomials}
\epsilon_A \Pi_{I \in A} P_I - \epsilon_B \Pi_{J \in B} P_J
\end{equation}
if and only if the contents of the corresponding $\Lambda$-tableau of size $k \times |A|$ are row-wise equal.
\end{Definition}

To $I \in \mathbf{I}_{k, n}$ {with $\sigma=\Lambda(I)$} we associate the monomial 
\[ 
\mathbf{x}_{\Lambda(I)}\coloneqq x_{\sigma(1) i_{1}}x_{\sigma(2)i_2}\cdots x_{{\sigma(k)i_k}}. 
 \]
A $k\times n$ matching field $\Lambda$, 
gives a map of polynomial rings 
\begin{equation}\label{eqn:monomialmap}
\begin{aligned}
\phi_{\Lambda} \colon\ \mathbb{K}[P_I] &\rightarrow \mathbb{K}[x_{ij}] 
\\ 
 P_{I} &\mapsto \sgn(\Lambda(I)) \mathbf{x}_{\Lambda(I)}.
\end{aligned}
\end{equation}

\smallskip



\begin{Proposition}
Given a matching field $\Lambda$, the matching field ideal $\mathcal{I}_{\Lambda}$ is the kernel of the monomial map
$\phi_{\Lambda}$ from Equation~\eqref{eqn:monomialmap}. 
\end{Proposition}


\setcounter{section}{3}
\section{Block diagonal matching fields}\label{sec:block}

In this section we describe a family of coherent matching fields of size $3 \times n$ whose toric ideals are generated in degree two and therefore yield toric degenerations and SAGBI bases of $\Gr(3,n)$. 



Consider a sequence of positive numbers $ a_1, a_2, \ldots, a_{r}$ so that $\sum_{i = 1}^r a_i = n$. For $1 \leq s\leq r$ set 
$I_s = \{\alpha_{s-1}+1, \alpha_{s-1} +2, \ldots, \alpha_{s}\}$,
where $\alpha_{s} = \sum_{i = 1}^{s} a_{i }$ and $\alpha_0 = 0$. 

\begin{Definition}\label{def:blockdiagcolumns}

The block diagonal matching field of size $3 \times n$ corresponding to a collection $\mathbf{a} = \{a_1, \ldots, a_{r}\}$ satisfying $\sum_{i = 1}^r a_i = n$ is denoted $\BLambda_{\mathbf{a}}$. This matching field is defined by:
\begin{enumerate}

\item $\BLambda_{\mathbf{a}}(I) = \id$ if $|I \cap I_s| \geq 2$ where $s$ is the minimal $t$ such that $I_t \cap I \neq \emptyset$;

\item \label{transpose} $\BLambda_{\mathbf{a}}(I) = (12)$ if $|I \cap I_s| = 1$ where $s$ is the minimal $t$ such that $I_t \cap I \neq \emptyset$. 

 \end{enumerate}
 A $2$-block diagonal matching field is a block diagonal matching field with $r = 2$. 
\end{Definition}

\begin{exam}
Consider the case when $a_1 = 1 $ and $a_2 = n-1$. Then $I_1 = \{1\}$ and $I_2 = \{2, \ldots , n\}$. 
Then $\BLambda_{1, n-1}(I) = \id$ if and only if $I \subset I_2$. Otherwise, we have $1 \in I$ and $1$ appears in the second row of column tableau corresponding to $I$.
The matching field $\BLambda_{1,n-1}$ is isomorphic to a pointed matching field $\Lambda$. This isomorphism is given by acting on $[n]$ by the transposition $(12)$. 
In fact, the $\Lambda$-tableaux are then the PBW-tableaux from~\cite{Feigen}.
 \end{exam}



 
\setcounter{theo}{3}


 
 In general there is a $\mathbb{Z}^4$-grading given by the number of elements of type $I_1$ in different rows of a tableau. A $3\times d$ tableau $T$ is of degree $(\alpha, \beta,\gamma, d-\alpha-\beta-\gamma)$ 
 where 
\begin{enumerate}
\item $\alpha=|\{\text{Content of row 3 of } T\}\cap I_1| = |\{I\in T: \ |I \cap I_1| = 3\}|$ 
\item $\beta=|\{\text{Content of row 1 of } T\}\cap I_1|-\alpha = |\{I\in T: \ |I \cap I_1| = 2\}|$ 
\item $\gamma=|\{\text{Content of row 2 of } T\}\cap I_1|-\alpha-\beta = |\{I\in T: \ |I \cap I_1| = 1\}|$ 
\end{enumerate}
For two $\Lambda$-tableaux $T,T'$ which are row-wise equal, these numbers are equal. This implies the following lemma.
 
% \medskip
 
\begin{lemm}\label{lem:homo}
The ideal of a block diagonal matching field
has a $\mathbb{Z}^4$-grading given by $(\alpha, \beta,\gamma, d-\alpha-\beta-\gamma)$
from above. 
 \end{lemm}
 
 
 

\begin{proof}[Proof of Theorem~\ref{Theorem:2block}]

Consider a binomial relation obtained from two $\Lambda$-tableaux $T$ and $T'$ of size $3 \times d$ where $d >2$ whose contents are row-wise equal. By applying quadratic changes to the tableaux (changes involving only two columns) we will reduce the degree of this relation thus proving that the matching field ideal is quadratically generated. 

Given a $\Lambda$-tableau $T$, arrange the columns so that the first columns are those for which the matching field assigns the identity permutation and to the last column the matching field assigns the transposition $(12)$. Let $C$ denote the subtableau formed by the first columns and let $D$ denote the subtableau formed by the last columns. 

The tableaux $C$ and $D$ can each be put into semi-standard format. 
In other words, we can rearrange both $C$ and $D$ so that all rows are in weakly increasing order, and the columns of $C$ are strictly increasing, whereas the columns of $D$ are arranged so that the first and second entries are permuted from the diagonal order. 

Now given a binomial relation obtained from two $\Lambda$-tableaux $T$ and $T'$. We assume that $T$ (respectively $T'$) is organized as a pair of subtableaux $C$, $D$ (respectively $C'$, $D'$) satisfying the requirements described above. 
 If the first columns of $T$ and $T'$ are equal, then we can cancel them from the binomial relation and it is not a minimal generator. 
 We let $I$ and $I'$ denote the first columns of $T$ and $T'$, respectively. 
Otherwise by Lemma~\ref{lem:homo} the matching field relations are homogeneous with respect to the $\mathbb{Z}^4$-grading and so $\lvert I \cap I_1 \rvert = \lvert I' \cap I_1\rvert$.

\let\oldthecdrthm\thecdrthm
\setcounter{cdrthm}{0}
\def\thecdrthm{\arabic{cdrthm}}
\begin{enonce}[remark]{Case}\label{Case1}%\refstepcounter{cdrthm}
Suppose that $\lvert I \cap I_1 \rvert = \lvert I' \cap I_1\rvert = 3$. In this case, the columns $I$ and $I'$ could only differ in the second row. Suppose the entries of the second row of $I$ and $I'$ are $j $ and $j'$, respectively. We can also assume that $j <j'$. 
Then there must be a $j$ in the second row of the tableaux $D'$ since the contents are row-wise equal and $C'$ is in weakly increasing order. Then swap the positions of $j$ and $j'$ in the second row of $T'$ so that the first columns of $T$ and $T'$ now agree. Notice that we can exchange the position of $j$ with that of $j'$ since $j, j' \in I_1$ and $j'$ was originally in the second row of a column whose the first and third entries were in $I_2$. 
\end{enonce}

\begin{enonce}[remark]{Case}\label{Case2}%\refstepcounter{cdrthm}
Suppose that $| I \cap I_1 | = |I' \cap I_1| =2$. In this case, the columns $I$ and $I'$ may only differ in the second and third rows but not in the first. Assume the column $I$ is $i, j, r$ and the column $I'$ is $i, j', r'$. If $j <j'$ then just as above there must be a $j$ in the second row of $D'$. We have that $j <r'$ since $j$ is in the first block and $r'$ is in the second block, so we can swap the positions of $j $ and $j'$ in the second row of $T'$. 

Assume now that $j = j'$, without loss of generality we can suppose that $r<r'$. Then there is an $r$ in the third row of $D'$. Suppose that the column containing $r$ is $s, t, r$. Then we can swap the positions of $r$ and $r'$ since $t<r<r'$ and we can place $r$ in the last row. Now the two first terms are equal and hence, the binomial is not a minimal generator.
\end{enonce}

\begin{enonce}[remark]{Case}\label{Case3}%\refstepcounter{cdrthm} 
Suppose that $| I \cap I_1 | = |I' \cap I_1| = 1$. In this case, the tableaux $C$ and $C'$ are empty. Then the first two columns must be equal since $T = D$ and $T'= D'$ and they are both in (transposed) semi-standard form. 
\end{enonce}

\begin{enonce}[remark]{Case}\label{Case4}\setcounter{cdrthm}{4}
Suppose that $| I \cap I_1 | = |I' \cap I_1| =0$. In this case, the entries of $I$ and $I'$ can only differ in the first and third row. Suppose that the column $I$ is $r, s, t$ and that the column $I'$ is $r', s, t'$. 
Therefore $r, r'<s<t, t'$ and we can assume that $r' <r$. 
Then there is a column in $T'$ with $r$ in the first row and we can swap $r$ and $r'$ in the first row of $T'$. 
Thus we may assume that $r = r'$ and without loss of generality that $t < t'$. Then there must be a $t$ somewhere in the last row of $D'$ and we can again swap $t$ and $t'$ so that the columns are now equal. This completes the proof. \qedhere
\end{enonce}
\end{proof}


\begin{exam} \label{example:thm2blockillustration}
Here we provide examples of the manipulations of the matching field tableaux used throughout the proof of Theorem~\ref{Theorem:2block}. 
The examples we consider are in $\Gr(3,8)$ with the block matching field given by $I_1 = \{1,2,3,4 \}$ and $I_2 = \{5,6,7,8\}$.


The following is an example of a tableau in the semi-standard form as described in the beginning of the proof of Theorem~\ref{Theorem:2block}, 
\[
T = 
\begin{tabular}{ccccc}
 1 & 2 & 5 & 5 & 6 \\
 3 & 3 & 6 & 2 & 4 \\
 4 & 7 & 8 & 7 & 8 \\
 \multicolumn{3}{c}{\upbracefill} & \multicolumn{2}{c}{\upbracefill}\\
 \multicolumn{3}{c}{$C$} & \multicolumn{2}{c}{$D$}\\
\end{tabular}.
\]
The three columns on the left belong to $C$, since the entries are in ascending order and the two columns on the right belong to $D$. The degree of $T$ with respect to the grading described in Lemma~\ref{lem:homo} $(1,1,2,1)$ which can be seen by calculating $\alpha = |\{(1,3,4)\}|$, $\beta = |\{(2,3,7)\}|$ and $\gamma = \{(5,2,7), (6,4,8)\}$. 

%\medskip

As an example of Case~\ref{Case1} in the proof, consider the tableaux
\[
T = 
\begin{tabular}{ccc}
 1 & 3 & 5 \\
 2 & 4 & 3 \\
 4 & 7 & 6 \\
 \multicolumn{1}{c}{\upbracefill} & \multicolumn{2}{c}{\upbracefill}\\
 \multicolumn{1}{c}{$C$} & \multicolumn{2}{c}{$D$}\\
\end{tabular}, 
\quad
T' = 
\begin{tabular}{ccc}
 1 & 3 & 5 \\
 3 & 4 & 2 \\
 4 & 6 & 7 \\
 \multicolumn{1}{c}{\upbracefill} & \multicolumn{2}{c}{\upbracefill}\\
 \multicolumn{1}{c}{$C$} & \multicolumn{2}{c}{$D$}\\
\end{tabular}.
\]
Then we may swap the entries $2$ and $3$ in the second row of $T'$ to obtain a row-wise equal tableau whose left most column is $(1,2,4)$. 

%\medskip
As an example of a tableau manipulation performed in Case~\ref{Case2} of the theorem consider the tableaux, 
\[
T = 
\begin{tabular}{ccc}
 1 & 5 & 6 \\
 2 & 1 & 3 \\
 8 & 6 & 7 \\
 \multicolumn{1}{c}{\upbracefill} & \multicolumn{2}{c}{\upbracefill}\\
 \multicolumn{1}{c}{$C$} & \multicolumn{2}{c}{$D$}\\
\end{tabular},
\quad
T' = 
\begin{tabular}{ccc}
 1 & 5 & 6 \\
 3 & 1 & 2 \\
 6 & 7 & 8 \\
 \multicolumn{1}{c}{\upbracefill} & \multicolumn{2}{c}{\upbracefill}\\
 \multicolumn{1}{c}{$C$} & \multicolumn{2}{c}{$D$}\\
\end{tabular}.
\]
First we can swap entries $2$ and $3$ in the second row of $T'$ to obtain tableau whose left most column is $(1,2,6)$. Then we swap entries $8$ and $6$ in the third row of $T$ to obtain a tableau whose left most column is also $(1,2,6)$. Note that we cannot swap entries $6$ and $8$ in the third row of $T'$.

\medskip
We omit examples of Case~\ref{Case3} in the proof since no swaps occur in this case. For an example of the manipulations made in Case~\ref{Case4}, we consider instead the block matching field 
$I_1 = \{1, 2\}$ and $I_2 = \{3,4,5,6,7,8 \}$. Consider then the two tableaux, 
\[
T = 
\begin{tabular}{cccc}
 4 & 5 & 3 & 6 \\
 5 & 7 & 1 & 2 \\
 6 & 8 & 5 & 7 \\
 \multicolumn{2}{c}{\upbracefill} & \multicolumn{2}{c}{\upbracefill}\\
 \multicolumn{2}{c}{$C$} & \multicolumn{2}{c}{$D$}\\
\end{tabular},
\quad
T' = 
\begin{tabular}{cccc}
 3 & 6 & 4 & 5 \\
 5 & 7 & 1 & 2 \\
 7 & 8 & 5 & 6 \\
 \multicolumn{2}{c}{\upbracefill} & \multicolumn{2}{c}{\upbracefill}\\
 \multicolumn{2}{c}{$C$} & \multicolumn{2}{c}{$D$}\\
\end{tabular}.
\]
First we swap entries $3$ and $4$ in the first row of $T'$ to obtain a tableau whose left most column is $(4,5,7)$. Then we may swap entries $6$ and $7$ in the last row of $T'$ to obtain a tableau whose left most column is $(4,5,6)$ which is identical to the left most column of $T$. Note that we can perform the swap of entries $3$ and $4$ in $T$ however we cannot swap entries $6$ and $7$ in $T$. 

\end{exam}
\begin{thebibliography}{27}
\bibitem[9]{Feigen} Evgeny Feigin, \emph{$\mathbb{G}_a^m$ degeneration of flag varieties}, Sel. Math., New Ser. 18 (2012), no. 3, 513--537.
\bibitem[27]{sturmfels1993maximal} Bernd Sturmfels and Andrei Zelevinsky, \emph{Maximal minors and their leading terms}, Adv. Math. 98 (1993), no. 1, 65--112.
\end{thebibliography}
\end{document}
